\chapter{Las potencias}\label{ch:g8:powers}

Dobla una hoja de papel por la \omterm{def:g3:division:half}{mitad} $10$ veces (¡si puedes!): su grosor
se duplica cada vez, y $10$ duplicaciones lo multiplican por
$2^{10} = 1024$. Las potencias son la abreviatura de la multiplicación repetida;
este capítulo monta la notación y sus reglas, con un papel especial
para las potencias de $10$.

\section{Definición}

\begin{definition}[Potencia]\label{def:g8:powers:def}
Para un número $a$ y un número entero $n \geq 1$:
\[
a^n = \underbrace{a \times a \times \dots \times a}_{n
\text{ factores}},
\]
que se lee «$a$ elevado a $n$»; $a$ es la \emph{base} y $n$, el
\emph{exponente}\index{exponente}. Nombres especiales: $a^2$ es «$a$ al
cuadrado» y $a^3$, «$a$ al cubo». Por convenio:
\[
a^1 = a,
\qquad
a^0 = 1 \ (a \neq 0),
\qquad
a^{-n} = \frac{1}{a^n}.
\]
\end{definition}

\begin{example}\label{ex:g8:powers:def}
$3^4 = 3 \times 3 \times 3 \times 3 = 81$; \quad
$10^3 = 1000$; \quad $5^{-2} = \frac{1}{25} = 0.04$; \quad
$(-2)^3 = -8$ y $(-2)^4 = +16$
(\cref{ex:g8:negprod:powers}). Cuidado: $a^n$ \emph{no} es
$a \times n$: \ $2^5 = 32$, no $10$.
\end{example}

\begin{omfigure}
\begin{tikzpicture}[scale=0.9]
  \foreach \n/\x in {0/0, 1/1.4, 2/2.8, 3/4.2, 4/5.6} {
    \pgfmathsetmacro{\h}{0.28*pow(2,\n)}
    \draw[omDef, fill=omDef!20] (\x,0) rectangle (\x+0.9,\h);
    \node[below, font=\small] at (\x+0.45,0) {$2^\n$};
    \pgfmathtruncatemacro{\v}{pow(2,\n)}
    \node[above, font=\small] at (\x+0.45,\h) {$\v$};
  }
\end{tikzpicture}

{\small Potencias de $2$: cada barra es el doble de la anterior. El crecimiento por
multiplicación repetida se dispara mucho más deprisa que el crecimiento por
\omterm{def:g1:addition:def}{suma} repetida.}
\end{omfigure}

\section{Las reglas de los exponentes}

\begin{theorem}[Reglas de los exponentes]\label{thm:g8:powers:rules}
Para una base $a$ no nula y \omterm{def:g8:powers:def}{exponentes} enteros $m$ y $n$:
\[
a^m \times a^n = a^{m+n},
\qquad
\frac{a^m}{a^n} = a^{m-n},
\qquad
\left(a^m\right)^n = a^{m \times n},
\qquad
(ab)^n = a^n b^n .
\]
\end{theorem}

\begin{proof}[Demostración contando factores]
$a^m \times a^n$ pone en fila $m$ factores $a$ seguidos de $n$ más:
$m + n$ en total. $\left(a^m\right)^n$ repite un bloque de $m$ factores
$n$ veces: $mn$ factores. $(ab)^n$ contiene $n$ letras $a$ y $n$
letras $b$, que se pueden reagrupar. La regla del \omterm{def:g3:division:remainder}{cociente} sale de
cancelar $n$ de los $m$ factores; con los convenios
$a^0 = 1$ y $a^{-n} = \frac{1}{a^n}$, sigue siendo cierta incluso cuando
$n \geq m$.
\end{proof}

\begin{example}\label{ex:g8:powers:rules}
\[
2^3 \times 2^5 = 2^8 = 256,
\qquad
\frac{7^6}{7^4} = 7^2 = 49,
\qquad
\left(5^2\right)^3 = 5^6,
\qquad
2^4 \times 5^4 = 10^4 .
\]
Una trampa: las reglas se aplican a una \emph{base común} (o a un
\omterm{def:g8:powers:def}{exponente} común, en la última). Ninguna regla simplifica $2^3 \times 5^2$ ---
hay que calcular: $8 \times 25 = 200$.
\end{example}

\section{Las potencias de diez}

\begin{proposition}[Potencias de diez]\label{prop:g8:powers:ten}
Para $n \geq 1$: $10^n = 1\underbrace{0\dots0}_{n}$ y
$10^{-n} = 0.\underbrace{0\dots0}_{n-1}1$. Las reglas de los \omterm{def:g8:powers:def}{exponentes}
dicen: multiplicar potencias de diez \omterm{def:g1:addition:def}{suma} los \omterm{def:g8:powers:def}{exponentes}.
\end{proposition}

\begin{example}\label{ex:g8:powers:ten}
$10^4 \times 10^3 = 10^7$; \quad
$\dfrac{10^2}{10^5} = 10^{-3} = 0.001$. Las cantidades grandes y pequeñas
se vuelven legibles:
\[
\text{mil millones} = 10^9,
\qquad
\text{una millonésima} = 10^{-6}.
\]
Combinadas con decimales: $3.2 \times 10^5 = 320\,000$ y
$4.7 \times 10^{-3} = 0.0047$ (desplaza el \omterm{def:g4:decimals:point}{punto decimal},
\cref{prop:g6:decimals:shift}). El uso sistemático de esta escritura ---
la \emph{notación científica} --- se desarrolla en el
\cref{ch:g9:fractions}.
\end{example}

\begin{example}[Órdenes de magnitud]\label{ex:g8:powers:magnitude}
La luz viaja a unos $3 \times 10^8$ m/s; un año tiene unos
$3.2 \times 10^7$ segundos. Un año luz es, por tanto, de unos
\[
3 \times 3.2 \times 10^{8+7} \approx 10 \times 10^{15} = 10^{16}
\text{ m}
\]
--- diez millones de miles de millones de metros. Las potencias de diez hacen que los cálculos
astronómicos quepan en una línea.
\end{example}

\begin{method}[Simplificar una expresión con potencias]\label{met:g8:powers:simplify}
\begin{enumerate}
  \item agrupa los factores base por base;
  \item aplica las reglas de los \omterm{def:g8:powers:def}{exponentes} dentro de cada base;
  \item calcula las potencias pequeñas que queden, o deja la respuesta como
        potencia si es grande.
\end{enumerate}
\end{method}

\begin{example}\label{ex:g8:powers:simplify}
\[
\frac{3^5 \times 3^{-2} \times 4^2}{3^2}
= \frac{3^{5 + (-2)}}{3^2} \times 16
= 3^{3 - 2} \times 16
= 3 \times 16 = 48 .
\]
\end{example}

\section{Ejercicios}

\begin{exercise}[$\star$]\label{exo:g8:powers:1}
Calcula:
\[
2^6, \qquad 3^3, \qquad 10^5, \qquad 1^{100}, \qquad 0.1^2, \qquad
6^0 .
\]
\end{exercise}

\begin{exercise}[$\star$]\label{exo:g8:powers:2}
Calcula:
\[
(-3)^2, \qquad -3^2, \qquad (-1)^{15}, \qquad (-5)^3, \qquad
\left(\tfrac{2}{3}\right)^2 .
\]
\end{exercise}

\begin{exercise}[$\star$]\label{exo:g8:powers:3}
Escribe como una sola potencia:
\[
7^4 \times 7^5, \qquad
\frac{2^9}{2^3}, \qquad
\left(10^3\right)^4, \qquad
5^6 \times 5^{-2}, \qquad
3^4 \times 7^4 .
\]
\end{exercise}

\begin{exercise}[$\star$]\label{exo:g8:powers:4}
Escribe como \omterm{def:g6:decimals:places}{número decimal}: $10^{-2}$; \quad $4 \times 10^3$; \quad
$2.5 \times 10^{-4}$; \quad $10^0$.
\end{exercise}

\begin{exercise}[$\star$]\label{exo:g8:powers:5}
Escribe con una potencia de diez: cien mil; una décima; diez mil
millones; $0.000\,001$.
\end{exercise}

\begin{exercise}[$\star$]\label{exo:g8:powers:6}
Simplifica y calcula después:
\[
\frac{10^7 \times 10^{-3}}{10^2},
\qquad
\frac{2^5 \times 2^4}{2^6},
\qquad
\left(2^2\right)^3 \times 2^{-4} .
\]
\end{exercise}

\begin{exercise}[$\star$]\label{exo:g8:powers:7}
¿Verdadero o falso? Corrige las falsas.
\[
2^3 \times 2^4 = 2^{12}; \qquad
5^2 + 5^3 = 5^5; \qquad
(3^2)^4 = 3^8; \qquad
10^3 \times 10^3 = 100^3 .
\]
\end{exercise}

\begin{exercise}[$\star\star$]\label{exo:g8:powers:8}
Un rumor se propaga: el día 1 lo conocen tres personas; cada día, cada
persona que lo conoce se lo cuenta a tres personas nuevas. Escribe con una potencia el número
de personas \emph{nuevas} informadas el día $4$ y calcula cuántas personas
conocen el rumor al final del día 4 (incluidas las tres iniciales).
\end{exercise}

\begin{exercise}[$\star\star$]\label{exo:g8:powers:9}
Una hoja de papel tiene $0.1$ mm de grosor, es decir, $10^{-4}$ m. Doblarla
duplica su grosor cada vez.
\begin{enumerate}
  \item expresa el grosor tras $10$ dobleces como un \omterm{def:g2:mult:def}{producto} y
        calcúlalo en centímetros ($2^{10} = 1024$);
  \item tras $42$ dobleces el grosor sería
        $2^{42} \times 10^{-4}$ m, con $2^{42} \approx 4.4 \times
        10^{12}$. Muestra que eso supera la distancia Tierra--Luna,
        de unos $3.8 \times 10^8$ m.
\end{enumerate}
\end{exercise}

\begin{exercise}[$\star\star$]\label{exo:g8:powers:10}
Ordena de menor a mayor, sin calculadora:
\[
2^{10}, \qquad 10^3, \qquad 3^6, \qquad 5^4 .
\]
(Calcula cada uno; $2^{10}$ y $10^3$ son vecinos famosos.)
\end{exercise}

\begin{exercise}[$\star\star\star$]\label{exo:g8:powers:11}
¿Cuál es mayor, $2^{100}$ o $10^{30}$? Usa
$2^{10} = 1024 > 10^3$ para comparar
$2^{100} = \left(2^{10}\right)^{10}$ con
$\left(10^3\right)^{10}$.
\end{exercise}

\section{Problema: El tablero de ajedrez y las potencias de dos}

\begin{problem}[{Problema de fin de semana --- la \omterm{def:g1:addition:def}{suma} geométrica
$1 + 2 + 4 + \dots + 2^{n-1} = 2^n - 1$, de una leyenda famosa a
los números binarios}]
\label{pb:g8:powers:1}
La leyenda: como recompensa por inventar el ajedrez, el sabio Sissa pidió
a su rey un grano de trigo en la primera casilla del tablero,
dos en la segunda, cuatro en la tercera --- duplicando de casilla en
casilla, hasta la sexagésima cuarta. El rey se rio de tanta modestia.
Este problema calcula lo que prometió el rey usando las reglas de los \omterm{def:g8:powers:def}{exponentes}
del \cref{thm:g8:powers:rules}, y termina donde lleva en secreto la
historia: los números binarios que hay dentro de todo ordenador.

\medskip
\noindent\textbf{Parte I --- El truco de la duplicación.}
Para $n \geq 1$, sea $S_n$ el número total de granos de las
primeras $n$ casillas.
\begin{enumerate}
  \item expresa el número de granos de la casilla $k$ como potencia de
        $2$. ¿Qué potencia hay en la casilla $64$?
  \item calcula $S_1$, $S_2$, $S_3$, $S_4$ y $S_5$, y compara
        cada uno con una potencia de $2$ cercana. Conjetura una fórmula para
        $S_n$.
  \item el \emph{truco de la duplicación}: escribe las \omterm{def:g1:addition:def}{sumas} $S_n$ y
        $2 \times S_n$ una debajo de la otra, resta y demuestra
        tu conjetura:
        \[
        S_n = 1 + 2 + 4 + \dots + 2^{n-1} = 2^n - 1 .
        \]
  \item ¿cuántos granos prometió el rey en total? Expresa la
        respuesta con una potencia de $2$ y completa la
        observación clásica: «el tablero entero contiene un grano menos que
        una sola casilla sexagésima quinta»;
  \item muestra que la segunda \omterm{def:g3:division:half}{mitad} del tablero (casillas $33$ a
        $64$) contiene \emph{exactamente} $2^{32}$ veces más granos
        que la primera \omterm{def:g3:division:half}{mitad}.
\end{enumerate}

\medskip
\noindent\textbf{Parte II --- ¿Cómo de grande es $2^{64}$?}
La comparación $2^{10} = 1024 > 10^3$ del
\cref{exo:g8:powers:11} es la clave de todas las estimaciones de abajo.
\begin{enumerate}[resume]
  \item muestra que
        $2^{64} = 2^4 \times \left(2^{10}\right)^6 > 1.6 \times
        10^{19}$.
  \item un grano de trigo pesa unos $0.05$ g, es decir,
        $5 \times 10^{-2}$ g. Muestra que el trigo prometido pesa
        más de $8 \times 10^{17}$ g y convierte eso en
        toneladas ($1$ tonelada $= 10^6$ g);
  \item el mundo entero cosecha actualmente unos $8 \times 10^8$
        toneladas de trigo al año. ¿Cuántos años de cosecha
        mundial prometió el rey, como mínimo?
  \item halla el menor número entero $n$ tal que
        $2^n > 10^6$ --- es decir, cuántas duplicaciones hacen falta para
        pasar de un millón. (Calcula $2^{19}$ y $2^{20}$ exactamente,
        usando $2^{10} = 1024$.)
  \item un tramposo te ofrece un sueldo mensual: $1$ céntimo el día
        $1$ y, cada día, el doble de la paga del día anterior. ¿Qué
        día supera por primera vez la paga \emph{diaria} sola el
        millón de euros ($10^8$ céntimos)? (Calcula $2^{26}$ y
        $2^{27}$ exactamente.)
\end{enumerate}

\medskip
\noindent\textbf{Parte III --- Pesas binarias.}
Una comerciante tiene cinco pesas: $1$, $2$, $4$, $8$ y $16$ gramos,
una de cada. Coloca algunas de ellas en un platillo de la balanza para
pesar mercancías en el otro.
\begin{enumerate}[resume]
  \item ¿qué pesas coloca para pesar $21$ g? ¿Y para pesar
        $27$ g?
  \item explica por qué cualquier objetivo de $16$ g o más \emph{tiene} que usar
        la pesa de $16$ g, y por qué cualquier objetivo de $15$ g o menos
        \emph{no} debe usarla. (La pregunta 3 te dice hasta dónde llegan
        como mucho las pesas $1, 2, 4, 8$.) Explica por qué el
        mismo razonamiento se repite con la pesa siguiente en tamaño, en
        cada etapa;
  \item deduce que todo objetivo entero de $1$ a $31$ g se puede
        pesar, y \emph{de una sola manera}: cada número
        entre $1$ y $31$ es \omterm{def:g1:addition:def}{suma} de potencias de $2$ distintas de
        una única forma;
  \item la comerciante compra una sexta pesa, de $32$ g. ¿Hasta qué
        objetivo puede pesar ahora? Escribe $45$ g como \omterm{def:g1:addition:def}{suma} de
        potencias de $2$ distintas;
  \item en un ordenador, un número «de 64 bits» se guarda en $64$
        casillas, cada una con un $0$ o un $1$ --- la casilla $k$
        aporta $2^{k-1}$ cuando lleva un $1$, como los
        granos de la leyenda. Usando la Parte I, explica por qué los números
        enteros que puede guardar una máquina así van exactamente de $0$ a
        $2^{64} - 1$.
\end{enumerate}
\end{problem}
